\documentclass[11pt]{article}

\usepackage[T1]{fontenc}
\usepackage[brazil]{babel}
\usepackage{indentfirst}
\usepackage{graphicx}
\usepackage{array}
\usepackage{geometry}
\usepackage{fancyhdr}
\usepackage{csquotes}
\usepackage{needspace}
\usepackage{float}
\usepackage{amssymb}
\usepackage{amsmath}
\usepackage{enumitem}
\setlist{nosep,leftmargin=*}
\usepackage{hyperref}
\usepackage{xurl}
\usepackage[style=numeric,sorting=nyt]{biblatex}
\addbibresource{referencias.bib}

\geometry{
    a4paper,
    width=160mm,
    top=20mm,
    bottom=40mm,
    bindingoffset=3mm
}

\pagestyle{fancy}
\fancyhead{}

\fancyhead[L]{%
    \parbox[c]{0.60\textwidth}{%
        \noindent{\large \textbf{NeuroLógica}}\\[2pt]
        {\small \textbf{Disciplina:} Redes Neurais I}\\[1pt]
        {\small \textbf{Professor:} Davi Pereira dos Santos}%
    }%
}

\fancyhead[R]{%
    \includegraphics[width=0.30\textwidth]{fig/uel_horizontal.jpg}%
}

\renewcommand{\headrulewidth}{0pt}
\setlength{\headheight}{58pt}
\setlength{\headsep}{1.8cm}
\fancyheadoffset{1cm}

\renewcommand{\topfraction}{0.9}
\renewcommand{\textfraction}{0.08}
\renewcommand{\floatpagefraction}{0.7}

\begin{document}

\begin{center}
    \Large\textbf{NeuroLógica}\\[2pt]
    \large Um neurônio traça uma reta: Perceptron, Adaline e MLP nas sete portas lógicas
\end{center}

\smallskip

\section*{Identificação}

\begin{itemize}
    \item \textbf{Integrantes:} Adriano Lúcio Uchôa Brandão, Herik Daurizio Ricardo, Heitor Oliveira Nunes
    \item \textbf{Disciplina:} Redes Neurais I, Universidade Estadual de Londrina (UEL)
    \item \textbf{Professor:} Davi Pereira dos Santos
    \item \textbf{Nome do projeto:} NeuroLógica
    \item \textbf{Questões orientadoras respondidas:} Q2, Q5 e Q4, nessa ordem no site
    \item \textbf{Repositório:} \url{https://github.com/cowpox/neurologica} (público)
    \item \textbf{Site:} \url{https://neurologica.vercel.app/}
\end{itemize}

\section{Explicação intuitiva}

Um neurônio artificial aprende traçando uma reta. Ele soma as entradas, cada uma
multiplicada por um peso, acrescenta um viés e decide: liga ou não liga. Essa soma define uma
reta. Treinar é mover essa reta, pouco a pouco, até ela separar as duas classes do problema.

O NeuroLógica usa as sete portas lógicas (AND, OR, NOT, NAND, NOR, XOR e XNOR) para contar
essa história em três atos, um por pergunta orientadora. No Ato 1, o Perceptron só corrige os
pesos quando erra. Ele resolve cinco das sete portas, mas trava no XOR e no XNOR: não existe
reta que separe esses dois casos (Q2). No Ato 2, o Adaline corrige os pesos sempre, mesmo
quando já acerta. Isso é o gradiente descendente em ação (Q5). No Ato 3, uma pequena rede
dobra o espaço até o XOR virar separável por uma única reta. Isso só funciona porque a
ativação de cada neurônio não é linear (Q4).

A ideia central, repetida nos três atos: a rede não resolve um problema difícil. Ela
transforma o problema difícil em um problema fácil, uma camada de cada vez.

\section{Referencial teórico}

\textbf{Do neurônio de McCulloch-Pitts ao Perceptron.} O neurônio de McCulloch e
Pitts~\cite{mcculloch1943logical} já somava entradas e decidia por um limiar, mas com pesos
fixos, escolhidos à mão: o NAND, por exemplo, exigia dois neurônios (AND seguido de NOT). O
Perceptron de Rosenblatt~\cite{rosenblatt1958perceptron} aprende os próprios pesos a partir de
exemplos:
$$z = \vec w \cdot \vec x + b, \qquad \hat y = \text{degrau}(z), \qquad w \leftarrow w + \eta\,(y-\hat y)\,x.$$
Com pesos aprendidos, um único neurônio já resolve o NAND. O algoritmo só corrige quando erra
a classe. Por isso ele para na primeira reta que zera o erro de treino, mesmo encostada num
ponto, como mostramos na seção de aprendizados.

\smallskip
\textbf{Adaline e o gradiente descendente (Q5).} O Adaline de Widrow e
Hoff~\cite{widrowhoff1960adaline} usa o erro de antes da ativação, $y - z$, não o erro de
depois:
$$w \leftarrow w + \eta\,(y - z)\,x.$$
Esse é um passo de gradiente descendente sobre o erro quadrático $E = \tfrac12\sum (y-z)^2$.
Por isso o Adaline é, na prática, uma regressão linear usada como classificador. Ele continua
ajustando os pesos mesmo depois de acertar os quatro pontos de uma porta, porque o erro
quadrático ainda não é zero.

\smallskip
\textbf{Por que uma reta não basta (Q2 e Q4).} Minsky e
Papert~\cite{minskypapert1969perceptrons} provaram que um único neurônio só resolve
problemas linearmente separáveis. O XOR não é um deles. A solução de Rumelhart, Hinton e
Williams~\cite{rumelhart1986backprop} é empilhar neurônios em camadas, com
\textit{backpropagation}. No site, XOR = AND(OR, NAND): os dois neurônios escondidos criam
um espaço novo onde uma última reta resolve tudo. Para isso funcionar, a ativação de cada
neurônio precisa ser não linear e diferenciável. Essa dupla exigência explica por que
sigmoide, tangente hiperbólica e ReLU existem lado a lado, cada uma com vantagens e problemas
próprios. O teorema da aproximação universal~\cite{cybenko1989approximation,hornik1991universal}
garante que a solução existe, mas não garante que o treino a encontra: a rede do Ato 3 tem
mínimos locais reais e precisa de reinícios automáticos.

\section{Conjunto de dados}

Não há um conjunto de dados externo. O dado do projeto são as sete portas lógicas
booleanas, cada uma com até quatro exemplos e rótulo conhecido pela própria tabela-verdade.
A escolha é proposital: é o menor problema visualizável em 2D, com solução exata conhecida
de antemão, e separa os casos que um único neurônio resolve dos que exigem uma rede. O
laboratório livre do site soma seis conjuntos sintéticos de classificação 2D (círculo, XOR
em quadrantes, duas gaussianas, espiral, meias-luas e xadrez), para quem quiser ir além das
sete portas.

\section{Disponibilização da implementação}

O código está publicado e é público, em \url{https://github.com/cowpox/neurologica}. O
repositório tem duas partes: \texttt{prototipos/}, os quatro marcos do projeto em Python, e
\texttt{docs/}, o site em HTML, CSS e JavaScript puro, publicado no Vercel em
\url{https://neurologica.vercel.app/}. A pasta \texttt{testes/} confere, com mais de 300
verificações automáticas, se o JavaScript do site reproduz os mesmos números dos protótipos
em Python.

\section{Aplicação e relevância atual}

Perceptron, Adaline e MLP com \textit{backpropagation} são a base conceitual de toda rede
profunda em uso hoje. O mesmo raciocínio do projeto (uma reta que separa classes, um passo
de gradiente que reduz um erro, camadas que dobram o espaço) é o raciocínio por trás de
redes muito maiores. A chamada \enquote{hipótese da variedade} torna essa ponte concreta:
dados reais de alta dimensão, como uma foto, não preenchem o espaço em que vivem. Eles se
concentram perto de superfícies de dimensão bem menor, e cada camada de uma rede pode ser
lida como uma deformação desse espaço até essas superfícies ficarem separáveis. É o que o
Ato 3 mostra em miniatura, com apenas quatro pontos.

\section{Estratégias para a apresentação ser compreendida pela turma}

A estratégia central é o próprio site: feito para a turma mexer ao vivo, via QR code, não só
para ser assistido. Cada ato propõe um desafio antes de qualquer explicação, por exemplo:
\enquote{treine o Perceptron nas sete portas e descubra quais ele resolve}. Conquistas e
pistas progressivas recompensam quem tenta sozinho antes de ler a resposta. A progressão
também é deliberada: o site mostra o Perceptron falhar e o Adaline quase resolver antes de
introduzir a rede, sempre com a mesma pergunta guia (\enquote{por que a reta pega esse ponto
e não aquele?}). Como plano B, se o site falhar na apresentação, os protótipos em Python
rodam localmente, com um modo \texttt{--comparar} que mostra o Perceptron e o Adaline lado a
lado.

\section{Histórico do desenvolvimento}

\begin{itemize}
    \item \textbf{04/09:} escolha do tema, uma animação de um neurônio ajustando retas para
    resolver portas lógicas.
    \item \textbf{06/09:} o projeto ganha nome (NeuroLógica), a pesquisa de trabalhos
    relacionados confirma o tema, e o grupo decide trabalhar em duas fases: primeiro Python,
    depois o site.
    \item \textbf{06--07/09:} Marco 1 (Perceptron), validado nas cinco portas separáveis.
    \item \textbf{09--12/10:} Marco 2 (Adaline); o Heitor entra no grupo, que fecha em três
    pessoas.
    \item \textbf{04/11:} apresentação intermediária pedida pelo professor.
    \item \textbf{10--20/11:} Marcos 3 e 4 (XOR e MLP) concluídos; o site inteiro é
    implementado e publicado no Vercel.
    \item \textbf{20--25/11:} revisão final, este documento e ensaio da apresentação.
\end{itemize}

\section{Relatos sobre a vivência de cada integrante}

\textbf{Adriano Lúcio Uchôa Brandão.} Minha parte foi tornar palpáveis conceitos abstratos:
por que o Perceptron para onde para, o que o gradiente descendente otimiza, o que muda numa
troca de função de ativação. Implementei o Marco 2 (Adaline) e participei da fundamentação
teórica. Também coordenei o grupo: dividi os marcos e as frentes do site entre os três,
acompanhei os prazos internos e revisei este documento. O que mais me marcou foi ver, nos
próprios experimentos, erros de arredondamento decidindo a convergência do treino, algo que
a teoria da disciplina não antecipa.

\smallskip
\textbf{Herik Daurizio Ricardo.} Minha parte foi implementar os algoritmos de aprendizado e
as animações da página. Depois dos quatro marcos validados em Python, portei toda a lógica
de treino para JavaScript, com testes que conferem se o site reproduz os mesmos números.
Isso revelou diferenças reais de arredondamento entre Python e JavaScript, que precisei
corrigir. Também fiquei responsável pela interface: os três atos, o laboratório livre e o
sistema de conquistas e pistas. Foi a parte em que mais aprendi a diferença entre uma
implementação que funciona em teoria e uma que funciona ao vivo, na tela de qualquer pessoa.

\smallskip
\textbf{Heitor Oliveira Nunes.} Minha parte foi cuidar da bibliografia do projeto e do
design inicial da página, além de ajudar a manter o grupo alinhado. Levantei os trabalhos
relacionados que situam o NeuroLógica diante do que já existe, como o TensorFlow Playground.
Contribuí com as primeiras decisões visuais do site, depois aprofundadas conforme a
interface ganhava forma. O maior aprendizado foi perceber que, num projeto deste tamanho, a
comunicação constante entre o grupo importa tanto quanto o código.

\section{Colaborações com outros grupos}

Houve troca de ideias com o grupo do colega Rafael Cobo, sobre animações de redes neurais e
ferramentas pedagógicas de aprendizado. Foi uma conversa natural, já que os dois projetos
exploram formas visuais e interativas de ensinar redes neurais, com temas e implementações
independentes. Não houve divisão de tarefas nem reaproveitamento de código entre os grupos.

\section{Aprendizados obtidos e sugestões para o futuro}

\textbf{O Perceptron não maximiza margem.} Em várias portas (OR, NOT, NAND, NOR), a reta
final termina encostada num ponto. Não é erro: o algoritmo só corrige quando erra, e para na
primeira reta que zera o erro de treino, sem noção de margem de segurança. Provamos que só
pontos da classe positiva podem tocar essa reta.

\smallskip
\textbf{Ruído de ponto flutuante pode decidir a convergência.} Com pesos iniciais zero, um
ponto do AND caiu repetidas vezes exatamente sobre a fronteira de decisão. Refeita a conta à
mão, com aritmética exata, a convergência real acontece duas épocas depois da que o
\texttt{float64} reportava. A correção foi adicionar uma tolerância numérica na comparação
com zero.

\smallskip
\textbf{O Adaline é uma regressão linear usada como classificador.} Isso explica por que ele
continua ajustando os pesos mesmo depois de acertar tudo, e por que, no XOR, o melhor ajuste
linear possível é a constante zero: um modelo que não decide nada. O limite está em o modelo
ser linear, não no gradiente descendente.

\smallskip
\textbf{Portabilidade numérica entre Python e JavaScript não é trivial.} Motores diferentes
de JavaScript não garantem o mesmo resultado, bit a bit, para funções como exponencial e
tangente hiperbólica. A solução foi reimplementar essas funções só com operações que o
padrão IEEE-754 garante idênticas em qualquer motor.

\smallskip
\textbf{Sugestões para o futuro.} Decidir mais cedo quem cuida de cada frente evitaria
realocações de tarefa no meio do projeto. Registrar cada decisão técnica num documento
próprio, no formato demanda, opções, escolha e motivo, ajudou a reconstituir o raciocínio
meses depois e vale recomendar a outros grupos.

\section{Ilustrações, links e referências}

\begin{itemize}
    \item \textbf{Site interativo:} \url{https://neurologica.vercel.app/}
    \item \textbf{Repositório público:} \url{https://github.com/cowpox/neurologica}
\end{itemize}

\smallskip

% TODO: inserir capturas de tela do site (Ato 1, Ato 2, Ato 3 e laboratório livre),
% geradas a partir do site publicado, em entrega/fig/, e descomentar os blocos abaixo.
%
% \begin{figure}[H]
%     \centering
%     \includegraphics[width=0.8\textwidth]{fig/screenshot_ato1.png}
%     \caption{Ato 1 do site: o Perceptron treinando na porta AND.}
%     \label{fig:ato1}
% \end{figure}

O projeto fica entre duas ferramentas parecidas. O TensorFlow
Playground~\cite{smilkov2017playground} é um sandbox genérico, sem portas lógicas nem
equação do hiperplano visível. O Perceptron Educational Lab~\cite{vlabs2016perceptron} é o
mais próximo do nosso escopo, mas cobre só três portas e não tem Adaline. O diferencial do
NeuroLógica é cobrir as sete portas do curso, ligar ao neurônio de McCulloch-Pitts, e mostrar
a progressão Perceptron, Adaline e MLP pelo gradiente descendente.

\section{Sugestões de aperfeiçoamento deste tipo de trabalho para turmas futuras}

A atividade intermediária ajudou a corrigir o escopo a tempo e vale manter em turmas
futuras. Por outro lado, a liberdade total de formato faz projetos grandes como o nosso
crescerem até perto do prazo. Um ponto de controle obrigatório, com um recorte mínimo já
funcionando, talvez ajude outros grupos a fechar o escopo mais cedo.

\section*{Declaração de uso de ferramentas generativas}

Este documento foi redigido com apoio do Claude Code (Anthropic). Os números citados foram
conferidos executando os scripts do projeto, não apenas gerados pela ferramenta. O Claude
Code também apoiou a implementação dos protótipos em Python e do site, sempre com
verificação manual e testes automatizados do grupo.

\clearpage
\printbibliography[title={Referências}]

\end{document}
